\subsection{All intrinsic second-commutator variations at the nine cusps}

\begin{statement}{E.82}{Proved}{The nine complete cusp-family obstruction maps, their kernels, and every third-stage fibre}
Retain the based plus Fermat cubic of E.36--E.37, with
$O=[1:-1:0]$, parameter $s=z/x$ and tangent one.
Let $N$ be its geometric free pro-three group of rank ten,
and use the actual local source and action of E.74--E.77.
Write
\[
 \begin{gathered}
 W_2=N/\gamma_3N,\qquad W_3=N/\gamma_4N,\qquad
 V=N^{\rm ab},\\
 A_2=\gamma_2N/\gamma_3N=\operatorname{Lie}_2V,
 \qquad A_3=\gamma_3N/\gamma_4N=\operatorname{Lie}_3V,\\
 T_r=\operatorname{Lie}_r(T_1),\qquad
 J_r=\operatorname{Lie}_r(J_1),\qquad
 T_1=A(1+3\zeta,7)
 \end{gathered}
 \tag{E82.1}
\]
in the integral basis of (E37.2)--(E37.4). Thus $\tau$
is the specified actual generator with elliptic block
$\left(\begin{smallmatrix}1&-3\\3&-2\end{smallmatrix}\right)$,
and $J_1$ is E.37's actual conjugation matrix.
The bases of $A_2,A_3$ are the integral Hall bases of
E.41, of ranks $45,330$.

For every cusp $i\in\{A_0,A_1,A_2,B_0,B_1,B_2,C_1,C_2,O\}$,
fix its unit tangent over $K=\mathbb Q(\zeta)$ and a
geometric path $p_i$ to it from the stated base point.
Let $\iota_i$ be its positive pro-three inertia generator,
and set
\[
 d_i=p_i\iota_i p_i^{-1},\qquad
 u_i=p_iT(p_i)^{-1}\in N,
 \qquad u_iT(d_i)u_i^{-1}=d_i^7.
 \tag{E82.2}
\]
The conventions for a path mean that
$p_i\iota_i p_i^{-1}$ is based at $O$.
The section families below are labelled by these fixed paths.
Their abelian inertia vectors in $V$ are
\[
 \bar d_i=e_i\quad(i\ne O),\qquad
 \bar d_O=-\sum_{i\ne O}e_i,
 \tag{E82.3}
\]
where the eight $e_i$ occur after $f_1,f_2$ in E.37's basis.

For $z\in A_2$, choose the continuous Hall representative
$w_z\in\gamma_2N$ described below. The pair
\[
           (u_iw_z,d_i)\quad\hbox{in }W_2^2
 \tag{E82.4}
\]
gives an actual local section over $\mathbb Q_3$ by E.77.
For fixed $i$ its complete geometric-conjugacy parameter
space is the compact abelian group
\[
 \begin{gathered}
 \mathcal P=A_2/(I-T_2)A_2^{J_2}
    =\operatorname{coker}G,
       \qquad G=(I-T_2)(I+J_2),\\
 \mathcal P\simeq
       \mathbb Z_3^{24}\oplus C_3^9\oplus C_9^5,
       \qquad C_m=\mathbb Z/m\mathbb Z.
 \end{gathered}
 \tag{E82.5}
\]
The nine resulting subsets of $H^1(\mathbb Q_3,W_2)$
are pairwise disjoint.

Let $\mathcal B_3$ be the actual local obstruction group
of E.74. On each entire family (E82.4), the obstruction
map is the following continuous homomorphism of compact groups:
\[
 \begin{gathered}
 \delta_i:\mathcal P\longrightarrow
 \mathcal B_3=
 \frac{A_3}{(T_3-7I)A_3+(J_3+I)A_3}
       \simeq C_3^{63}\oplus C_9^{23}\oplus C_{27}^7,\\
          \delta_i([z])=[7[z,\bar d_i]].
 \end{gathered}
 \tag{E82.6}
\]
In particular every value uses the full intrinsic
rank-$330$ central lattice. The image, its exact kernel,
and its cokernel are as follows:
\[
\begin{array}{c|c|c}
 i&\operatorname{im}\delta_i&\ker\delta_i\\ \hline
 A_0,B_0,O&C_3^{11}\oplus C_9^3&
     \mathbb Z_3^{24}\oplus C_3^3\oplus C_9^4\\
 A_1,A_2,B_1,B_2,C_1,C_2&
     C_3^{22}\oplus C_9^7\oplus C_{27}&
     \mathbb Z_3^{24}\oplus C_3^6
\end{array}
 \tag{E82.7}
\]
\[
 \begin{array}{c|c|c}
 i&\operatorname{coker}\delta_i&
       [\mathcal P:\ker\delta_i]\\ \hline
 A_0,B_0,O&C_3^{52}\oplus C_9^{20}\oplus C_{27}^7&3^{17}\\
 A_1,A_2,B_1,B_2,C_1,C_2&
      C_3^{41}\oplus C_9^{16}\oplus C_{27}^6&3^{39}
 \end{array}
 \tag{E82.8}
\]
The certificate supplies a separate matrix and both coordinate
inverses for each of the nine entries; equality of the displayed
invariants does not remove their marked maps.

The image of the third-stage section map over (E82.4) is
exactly $\ker\delta_i$. Every nonempty fibre, with a particular
lift chosen by (E82.13) below, is homeomorphic to
\[
 \mathbb Z_3^{330}\oplus C_3^7\oplus C_9^{56}
                    \oplus C_{27}^{25}\oplus C_{81}^3.
 \tag{E82.9}
\]
The zero parameter in each family has the actual full
$N$-section (E82.2); the parameters outside the stated
kernel have no $W_3$-section above them.
In particular the single Hall vector $z=[f_2,f_1]$
has nonzero obstruction at every one of the nine cusps;
its full $93$ coordinates and exact order are saved.
At $A_0$ its first obstruction coordinate is
$2\bmod3$. Thus the literal pair
$(u_{A_0}[q_2,q_1],d_{A_0})$, where $q_j$ is the
chosen lift of $f_j$, has no intrinsic third-stage lift.
The actual maps from $N$ to its intrinsic metabelian and
single-cusp marked quotients identify their first three
nilpotent stages with those used here. Thus their
third-stage lifting maps on these families have the same
images and fibres.

Here are exact coordinate maps for (E82.6)--(E82.9).
Let $M=[\,T_3-7I\mid J_3+I\,]$, and take the retained
matrices $L,R$ of E.74, with
\[
 LMR\equiv [D\mid0]\pmod{729},\qquad
 D=\operatorname{diag}(3^{v_1},\ldots,3^{v_{330}}).
 \tag{E82.10}
\]
The multiplicities of $v=0,1,2,3$ are $237,63,23,7$.
Both $L$ and $R$ have their recorded inverse modulo $729$,
and hence define invertible matrices over $\mathbb Z_3$.
Let $A_i$ be the $330$ by $45$ integral matrix
$z\mapsto[z,\bar d_i]$. Selecting the $93$ rows with
$v_j>0$ gives the actual finite homomorphism
\[
 \begin{gathered}
 F_i:A_2\longrightarrow
       \bigoplus_{v_j>0}\mathbb Z/3^{v_j},\qquad
 (F_iz)_j=(L\,7A_iz)_j\bmod3^{v_j},\\
 \Lambda_i=\ker F_i=K_i\mathbb Z_3^{45},\qquad
 \ker\delta_i=\Lambda_i/G\mathbb Z_3^{45}.
 \end{gathered}
 \tag{E82.11}
\]
Every $K_i$ is a saved integer Hermite basis, with exact
rational inverse, determinant $3^{17}$ or $3^{39}$ as in
(E82.8), and $27K_i^{-1}$ integral. The matrices
$K_i^{-1}G$ are integral. Exact Smith maps over
$\mathbb Z_{(3)}$ and both their inverses are retained for
$G$, all $K_i$, and all $K_i^{-1}G$.
In detail, write
\[
 \begin{gathered}
 L_iK_iR_i=D_i,\qquad
 P_i(K_i^{-1}G)Q_i=E_i.\\
 \operatorname{im}\delta_i:\quad
     z\longmapsto L_iz\pmod{D_i},\qquad
     \eta\longmapsto F_iL_i^{-1}\eta;\\
 \ker\delta_i:\quad
     z\longmapsto P_iK_i^{-1}z\pmod{E_i},\qquad
     \eta\longmapsto K_iP_i^{-1}\eta\pmod{G}.
 \end{gathered}
 \tag{E82.12}
\]
The first line of coordinates identifies an obstruction value
using any parameter preimage $z$; the saved generators
$F_iL_i^{-1}\eta$ give its inverse. For $\mathcal P$ itself
the corresponding coordinates use the saved Smith maps of $G$.
All nonzero Smith entries have been scaled to powers of three
by actual units in $\mathbb Z_{(3)}$.

For a particular geometric third-stage lift, put
$v=7A_iz$, take the first $330$ columns $R_*$ of $R$, and set
\[
 \begin{gathered}
 C=D^{-1}LMR_*,\qquad C\equiv I\pmod{27},\qquad
 C^{-1}=\sum_{n\ge0}(I-C)^n,\\
 \binom{d}{e}=-R_*C^{-1}D^{-1}Lv,\qquad
        (u_iw_z,\ d_i\,d)\quad\hbox{in }W_3^2.
 \end{gathered}
 \tag{E82.13}
\]
Here $z\in\Lambda_i$, and $d$ in the last pair denotes
the central group element with its displayed coordinate in
$A_3$; $e$ is an auxiliary relation coordinate. This is a
convergent formula in $\mathbb Z_3$, with no unevaluated
higher Frobenius coefficient. The certificate verifies the
actual corrections modulo $729$ on all $405$ columns of the
nine kernel bases.

Finally let $D_V\subset V$ be the eight-dimensional cusp
lattice generated by the $e_i$, and define the homomorphism
\[
 \begin{gathered}
 A_2^8\longrightarrow A_3,\qquad
       (z_i)_{i\ne O}\longmapsto\sum_{i\ne O}[z_i,e_i],\\
 \mathscr I=\sum_i\operatorname{im}\delta_i
       \subset\mathcal B_3.
 \end{gathered}
 \tag{E82.14}
\]
The first image is a saturated direct summand of rank $304$,
with a specified Hall-coordinate complement of rank $26$.
The sum in the second line is the image of the induced
homomorphism $\mathcal P^8\to\mathcal B_3$, and has exact
groups
\[
 \mathscr I\simeq C_3^{61}\oplus C_9^{25}\oplus C_{27}^2,
 \qquad
 \mathcal B_3/\mathscr I\simeq C_3^{11}\oplus C_9.
 \tag{E82.15}
\]
The quotient map in (E82.15), its finite coordinates and their
inverse are given by the saved $93$ by $93$ relation-lattice
Smith maps. Thus the nine families evaluate a specified
subgroup of the full obstruction group, with its remaining
quotient explicitly retained.

\emph{Manuscript: \S51.16, folios733v--736v,
scans172--178, formulas(371)--(379), Part E(c),(d)
and follow-up(4). These are full intrinsic local families
and all their third-stage fibres; the remaining global and
higher intrinsic equations stay in E.23.}
\end{statement}

\begin{proof}
\emph{Actual paths and sections.}
The rational coordinate permutation
$[x:y:z]\mapsto[z:y:x]$ carries the base tangent of
E.34 to the stated point $O$ and changes its parameter
by the rational factor $-1$. The compatible three-power
roots of this factor are all $-1$. This is the actual
Galois-invariant tangent comparison used in E.78, so the
based action still factors through the same actual
$P\rtimes\langle c\rangle$ of E.74.
We use its generator with the specified E.37 matrix in
(E82.1).

Every cusp and chosen unit tangent is defined over $K$.
On its local inertia the actual $\tau$ acts by the seventh
power. Applying $T$ to the path expression for $d_i$ and
cancelling $T(p_i)^{-1}T(p_i)$ gives (E82.2) as an identity
in the full group $N$. In particular the relator value
$u_iT(d_i)u_i^{-1}d_i^{-7}$ is one. E.77 now gives an
actual section over $\mathbb Q_3$, with values
$u_i,J(u_i),d_i,J(d_i)$ on its four specified generators
and trivial value on its relative order-two subgroup.
The latter subgroup lies in the actual relative source
of E.74, with its proved inflation map from the absolute
local group.

Choose lifts of the two elliptic basis vectors of $V$
and use the eight fixed $d_i$, $i\ne O$, for the other
eight generators. Their classes form a basis of $V/3V$.
They generate $N$ by the pro-three Frattini criterion.
The induced surjection from a free pro-three group of
rank ten to $N$ is an isomorphism, by the Hopfian property
used in E.12. They are therefore an actual free basis.
For $z=(z_b)$ in its $45$ Hall coordinates define $w_z$
to be the ordered product of the corresponding group
commutators raised to their $\mathbb Z_3$ powers $z_b$.
This is a continuous map into $\gamma_2N$ whose image
modulo $\gamma_3N$ is exactly $z$.
The relation among all peripheral cocharacters gives
(E82.3). A change of the lift $w_z$ by $\gamma_3N$ has
no effect on (E82.4).

In $W_2$ the subgroup $A_2$ is central, so multiplication
of $u_i$ on the right by $w_z$ leaves its relator value
one. This proves the section assertion for every parameter,
not just for parameters in the later lifting kernel.

\emph{The entire parameter gauge.}
By E.77 two pairs (E82.4) are geometrically conjugate
precisely by its residual $J$-fixed geometric conjugators.
Let $g\in W_2^J$ carry the parameter $z$ to $z'$.
Comparing the abelian coordinates of their first entries
gives $(I-T_1)\bar g=0$. The elliptic block has
characteristic polynomial $X^2+X+7$, whose value at one
is nine, and the cusp block is $7I_8$.
Consequently $I-T_1$ is invertible over $\mathbb Q_3$,
and its kernel on the torsion-free integral lattice is
zero. Thus $g$ belongs to $A_2$ and is a vector
$b\in A_2^{J_2}$. The exact gauge formula of E.77 becomes
\[
                  z'-z=(I-T_2)b.
 \tag{E82.16}
\]
Conversely every such central $b$ fixes $d_i$ and gives
exactly this change. This proves the full orbit relation.
Since two is a unit in $\mathbb Z_3$, the image of
$I+J_2$ is precisely $A_2^{J_2}$: the inverse on the
fixed lattice is $b\mapsto b/2$. This proves (E82.5)
with the actual matrix $G$. Its Smith computation has
seven unit factors, nine factors three, five factors
nine and $24$ zero factors, in agreement with E.41.
Conjugation leaves the second entry's abelian vector
unchanged. The nine vectors (E82.3) are distinct, so
the nine families have disjoint geometric-conjugacy images.

\emph{The exact intrinsic third-degree error.}
Compute in $W_3$. Its third subgroup is central, and
$[w_z,T(d_i)]$ has coordinate $7[z,\bar d_i]$.
The group identity $waw^{-1}=[w,a]a$ and the full
identity (E82.2) give
\[
 \begin{aligned}
 (u_iw_z)T(d_i)(u_iw_z)^{-1}d_i^{-7}
   &=u_i[w_z,T(d_i)]u_i^{-1}
      u_iT(d_i)u_i^{-1}d_i^{-7}\\
   &=[w_z,T(d_i)]=7[z,\bar d_i]\quad\hbox{in }A_3.
 \end{aligned}
 \tag{E82.17}
\]
All higher terms of the original path and Frobenius
cancel in this equality. Changing $w_z$ by a third-degree
element leaves the displayed commutator unchanged.
Writing $v$ for (E82.17), E.77 gives the defect
$(I-J_3)v$ and the local invariant class $[(I-J_3)v/2]$.
In the quotient by $(J_3+I)A_3$ this class is $[v]$,
which proves (E82.6), with its factor seven and its sign.

The operator identity needed for the parameter quotient
is an identity of integral Lie maps:
\[
       T_3[z,\bar d_i]=[T_2z,7\bar d_i],\qquad
 7[(I-T_2)b,\bar d_i]=-(T_3-7I)[b,\bar d_i].
 \tag{E82.18}
\]
Here $T_1\bar d_i=7\bar d_i$ is the actual inertia
character. The second equality shows that every gauge
relation maps to zero in $\mathcal B_3$, and proves
that $\delta_i$ is a homomorphism on the full group
$\mathcal P$. The matrices $A_i$ and both identities
are checked integrally in the certificate.

\emph{All finite coordinates and their inverse maps.}
E.74 proves that the obstruction cokernel is killed by
$3^5$ before making its finite computation modulo $3^6$.
Its actual Smith calculation (E82.10) has maximum
valuation three. Therefore multiplication by $L$ and
reduction in its positive rows identifies the exact
completed cokernel with the displayed sum of $93$
cyclic groups. This gives the map $F_i$ in (E82.11).

We describe the complete kernel computation, including
the passage from its integer certificate to the compact
lattice. Start with $K=I_{45}$. At stage $j=0,1,2$ use
the rows of $F_i$ having exponent greater than $j$.
All entries of their product with the current $K$ are
divisible by $3^j$. Divide by that power and row-reduce
over $\mathbb F_3$. For each pivot column take its unit
vector; for each free column subtract the pivot-vector
combination prescribed by that reduced column. These
vectors form an integer unimodular matrix: a permutation
of its rows puts it in block triangular form with
identity diagonal blocks. Multiply its pivot columns
by three, multiply by $K$, and take the integer Hermite
basis of the resulting full lattice.

At this step a vector satisfies the next congruence
exactly when its pivot coordinates in that unimodular
basis are multiples of three. Induction therefore
proves equality with the solution lattice after each
stage, over both $\mathbb Z$ and $\mathbb Z_3$.
After three stages all $93$ congruences have been
imposed. The resulting matrix is the claimed $K_i$.
The successive pivot counts are $(14,3,0)$ at
$A_0,B_0,O$ and $(30,8,1)$ at the other six cusps.
They give the determinants in (E82.8).
The direct checks $F_iK_i=0$ in each target cyclic
group and $27K_i^{-1}$ integral also retain the two
containments and the exponent bound.

For completeness the exact Smith operation used on
these bases works in the localization $\mathbb Z_{(3)}$.
Choose an entry of minimum three-adic valuation in the
remaining block, move it to the next diagonal position,
and multiply its row by the inverse of its unit part.
The pivot is now exactly $3^a$. Every other entry in
that block is divisible by $3^a$ in $\mathbb Z_{(3)}$,
so row and column additions clear its column and row
with coefficients in that ring. Earlier pivots remain
unchanged. The remaining entries still have valuation
at least $a$, proving that the resulting diagonal is
a Smith diagonal. The program tracks the left and right
operations and their inverse operations, and separately
multiplies all full matrices to verify
$LAR=D$, $LL^{-1}=L^{-1}L=I$, and
$RR^{-1}=R^{-1}R=I$. Every denominator in these four
matrices is prime to three. Extending these exact maps
to $\mathbb Z_3$ proves the stated continuous coordinate
isomorphisms.

For $K_i$ the positive valuations are $(1^{11},2^3)$
or $(1^{22},2^7,3)$, respectively. For $K_i^{-1}G$
they are $(1^3,2^4)$ or $(1^6)$, with $24$ zero rows
in both cases. Their cokernels prove (E82.7) and
the explicit two-sided coordinate maps (E82.12).
The actual matrices $F_iL_i^{-1}$ are saved modulo
the target cyclic orders, using the inverse of every
rational unit denominator. Thus the obstruction
generators themselves, and not just their orders,
are retained.
The first degree-two Hall vector is $[f_2,f_1]$.
Its first matrix column is nonzero in each of the nine
actual cyclic targets, as recorded separately in the
certificate. At $A_0$ its first row is two modulo
three, proving the stated concrete nonlifting pair.

There is also an independent finite check of every
image invariant. For $k=0,1,2,3$, compute the cokernel
of the actual matrix
\[
 [\,\operatorname{diag}(3^{v_j})\mid 3^kF_i\,]
              \quad(v_j>0)
 \tag{E82.19}
\]
modulo $81$. These cokernels are killed by $27$
because the diagonal columns already impose that
bound. Minimum-valuation elimination modulo $81$
therefore gives their exact Smith valuations.
Subtracting their order exponents from $130$ gives
$\log_3|3^k\operatorname{im}F_i|$.
Successive differences recover the multiplicities
of cyclic factors of orders three, nine and twenty-seven.
They agree with the exact $K_i$ computations and also
give the cokernels in (E82.8).

\emph{Particular lifts and all full conjugacy fibres.}
Equation (E82.10), with its integer representatives,
implies $LMR_*=D+729E$ for an integer matrix $E$.
Since every diagonal entry of $D$ divides $27$,
$C=D^{-1}LMR_*$ is integral and congruent to the
identity modulo $27$. Its inverse series in
(E82.13) converges uniformly on the compact lattice.
For $z\in\Lambda_i$, the vector $D^{-1}Lv$ is
integral by (E82.11). Thus both $d,e$ in (E82.13)
belong to $A_3$, and multiplication gives the exact
identity
\[
              (T_3-7I)d+(J_3+I)e=-v.
 \tag{E82.20}
\]
Apply $I-J_3$ to this equality. The second summand
vanishes, and the first cancels the exact defect
$(I-J_3)v$ from (E82.17). E.77 therefore proves
that the displayed corrected pair is an actual
$W_3$-section. The stored correction matrix checks
this equality on every kernel basis column modulo
$729$. At that precision the inverse series is
$I-(C-I)$, since $(C-I)^2$ is divisible by $729$.
The convergent formula proves the lift at every
precision.

Conversely a third-stage section above (E82.4)
has zero local obstruction, so its parameter lies
in $\Lambda_i$. This establishes the entire image.
E.74's actual central-difference map now identifies
each nonempty full geometric-conjugacy fibre with
$H^1(\mathbb Q_3,A_3)$, giving (E82.9).
More explicitly one multiplies the particular lift
by a central cocycle; the inverse first aligns its
lower section by the unique lower geometric
conjugator and then takes the central difference.
The absence of a lower stabilizer is E.74's positive
weight argument for these actual modules. E.75 and
E.77 give the cocycle coordinate inverse
$A_3/(T_3-I)A_3^{J_3}\oplus
\ker((I-J_3)(T_3-7I))$.
Every map just used is continuous; the orbit spaces
are compact Hausdorff. Their specified continuous
inverse maps prove the homeomorphism assertions.

The quotient maps to the intrinsic metabelian and
marked groups are the actual maps of E.41. They
induce isomorphisms on the first three integral
graded groups. In the finite lower-central
filtration of $N/\gamma_4N$, a kernel element
therefore maps trivially in grade one, then in
grade two, then in grade three, and is the
identity. The maps onto the corresponding
third-stage groups are surjective by their
definition, hence are isomorphisms. They retain
the arithmetic action and the source markings,
which proves the stated transfer of section
images and fibres.

\emph{The sum of the nine evaluated images.}
The chosen Hall order first places the two elliptic
letters and then the eight cusp letters. A degree-three
Hall word is $[[a,b],c]$ with $a>b$ and $b\le c$.
Every such word whose outer letter $c$ is a cusp
letter is itself a column of (E82.14). The remaining
Hall words number $26$: there are two involving
only the elliptic letters and three for each of
the eight cusp letters. They are the words with
outer elliptic letter and with its allowed inner
smaller elliptic letter. A column $[[a,b],e_i]$ that
is not already Hall has $b>e_i$. Then $a,b,e_i$ are
all cusp letters, and its Jacobi rewrite has outer
letters among $a,b$, again cusp letters. This proves
that every column has zero coordinates in the $26$
complementary rows. Independent associative tensor
expansions in the program check this integral rewrite.
The other $304$ rows have distinct unit
columns. Selecting those columns is an actual
integer inverse on this summand, proving both
its saturation and the stated complement.

The $O$-map is the negative sum of the eight others
by (E82.3). Hence the subgroup $\mathscr I$ is
exactly the image of the saved concatenated matrix
$[F_i]_{i\ne O}$. Applying (E82.19) to this matrix
gives
\[
       \log_3|3^k\mathscr I|=117,29,2,0
                \quad(k=0,1,2,3).
 \tag{E82.21}
\]
These values give the first decomposition in
(E82.15), and its presented cokernel has eleven
factors three and one factor nine.
For actual quotient coordinates form the integer
column lattice of
$[\operatorname{diag}(3^{v_j})\mid F_{A_0}\mid\cdots\mid F_{C_2}]$
in $\mathbb Z^{93}$. Its saved Hermite basis has
determinant $3^{13}$. Every presentation column
lies in it; the finite index $3^{13}$ just computed
proves equality of the lattices. Its exact local
Smith maps and both inverses therefore give the
quotient coordinates claimed in (E82.15).

The programs \texttt{e\_cusp\_family\_obstruction.py}
and \texttt{e\_cusp\_local\_smith.py} reconstruct all
of these data. They check all $405$ Hall brackets
against independent associative tensor expansions,
the nine $T$ and nine $J$ intertwiners, all original
E.74 inverse matrices, every complete parameter
kernel and gauge map, the $405$ particular central
corrections, and the joint direct-summand and quotient
inverses. The saved JSON contains the entire matrices,
kernel derivations, exact rational Smith transformations
and inverse maps; it has no unspecified arithmetic
coefficient in the nine obstruction homomorphisms.
\end{proof}
